\section{Conclusion}

\textbf{Interpretation.} The SVM baseline (90.18\% accuracy) trails all four deep models by approximately 7 percentage points, with the gap most pronounced on minority classes (macro recall 84.11\% vs.\ weighted 90.18\%). End-to-end fine-tuning allows deep models to adapt representations to the target distribution and achieve nearly uniform per-class performance, whereas frozen features generalize poorly to underrepresented classes. All four deep architectures perform within 0.34 pp (96.86--97.20\%), suggesting that with sufficient pretraining, architecture choice has limited impact on Caltech-101. Our ablations show that AdamW outperforms SGD by 1.58 pp with faster convergence, and that removing data augmentation does not degrade validation accuracy (97.74\% vs.\ 97.51\%), as ImageNet-pretrained features already encode augmentation-invariant representations.

\textbf{Lessons learned.} Pretraining is essential: even frozen CNN features vastly outperform hand-crafted ones (e.g., HOG). Optimizer and learning rate matter---AdamW is preferable for fine-tuning, and SGD requires careful tuning to avoid severe underfitting. Augmentation is not always beneficial and should be validated on the target dataset. Classical methods are less robust to class imbalance than deep models.

\textbf{Conclusion.} This benchmark establishes a reproducible comparison of classical and deep learning methods on Caltech-101 under consistent protocols. The findings clarify the gains from end-to-end fine-tuning over fixed-feature classification and the limited marginal impact of architecture choice when pretraining is strong. Future work could extend this framework to larger or noisier datasets, where augmentation and architectural differences may play a more pronounced role.
